\section*{Methods}

\subsection*{Reproducibility experiments}
Seven classification datasets from the TabPFN benchmark suite were obtained from the public reproduction repository included in our project. For the baseline comparison, each dataset was split once into training and test partitions, stratified by class, with 15\% held out. TabPFN, CatBoost, XGBoost, LightGBM, random forest, logistic regression and AutoGluon were run with default settings. For the generation comparison, TabPFN v1, v2, v2.5, v2.6 and v3 and TabFM were evaluated with training sets capped at 1,024 rows by stratified subsampling. ROC AUC used one-vs-rest averaging for multiclass datasets. Runs that failed are reported as missing and not imputed.

\subsection*{Cancer cohorts and outcomes}
We used harmonized molecular and clinical data for BRCA, ESCA, HNSCC, LSCC and LUAD from the TCGA PanCancer Atlas and CPTAC proteogenomic studies.\citep{tcga2012brca,tcga2017esca,tcga2015hnsc,tcga2012lusc,tcga2014luad,liu2018clinical,li2023pancancer} Molecular features comprised RNA expression, copy number, DNA methylation and somatic mutation; clinical covariates were excluded from all models. For a horizon $h$ of 1,095 days (3 years) or 1,825 days (5 years), death on or before $h$ was labeled 1 and follow-up beyond $h$ was labeled 0; patients censored before $h$ were excluded. The extreme-survival endpoint compared death within 3 years with survival beyond 5 years.

\subsection*{Frozen evaluation}
Each task used five repeats of five patient-grouped outer folds, with about 20\% of patients in each test set and the remainder divided into training (about 64\% overall) and validation (about 16\%). All 400 test sets were frozen with checksums before any model was run. Pooled tasks were harmonized on exact shared feature identifiers, with features absent from a cohort set to zero. Within each fold, the 100 highest-variance molecular features were selected on training patients only, which avoids the selection bias that arises when features are chosen using all patients.\citep{ambroise2002selection} Validation patients were used only to choose a classification threshold.

\subsection*{Models}
We compared logistic regression, random forest, CatBoost, XGBoost, LightGBM, AutoGluon (180-second limit per fold), TabFM (PyTorch implementation) and TabPFN v2, v2.5, v2.6 and v3, all on CPU at default settings, with no hyperparameter tuning of any model, and seed 42.\citep{breiman2001random,chen2016xgboost,ke2017lightgbm,prokhorenkova2018catboost,erickson2020autogluon,hollmann2025accurate,google2026tabfm,tabfmsoftware2026} AutoGluon's 180-second limit is far shorter than the one- to four-hour budgets of its original benchmark.\citep{erickson2020autogluon} TabPFN-3 was run from the local checkpoint, and hosted-API variants were not evaluated.\citep{priorlabs2026tabpfn3} TabFM is documented in a Google announcement and a software repository. Model code was taken from the TabPFN repository at commit \texttt{c0bc41a}, the TabPFN-3 checkpoint repository at commit \texttt{2997265} and the TabFM repository at commit \texttt{5ee6cd7}. TabPFN v2 to v2.6 refuse CPU inference with more than 1,000 training samples by default, a runtime safeguard and not their recommended data size, so the pooled 3-year task (about 1,050 training patients) was not evaluated for those versions.
% TODO before submission: add exact package versions (pip freeze of the evaluation image and local environment).

\subsection*{Metrics and statistics}
ROC AUC was the primary measure; precision-recall AUC, log loss, Brier score and threshold-dependent metrics were also recorded. Precision-recall AUC depends on event prevalence and was read against it,\citep{saito2015precision} and a lower log loss or Brier score\citep{brier1950verification} indicates better aggregate probability estimates without establishing calibration.\citep{van2019calibration} We use two ROC AUC summaries and name them distinctly. Fold-mean AUC is the mean over the 25 test sets of a task (Tables~\ref{tab:model_auc} and~\ref{tab:shortcuts}, Figs.~\ref{fig:within} and~\ref{fig:pooling}a). Patient-averaged within-cancer AUC is computed after averaging each patient's test predictions across the five repeats (Fig.~\ref{fig:effect}). Models were compared across the 13 cancer-specific tasks with the Friedman test on per-task mean AUC and the Nemenyi critical difference for $k = 11$ models;\citep{demsar2006statistical} non-significance was not treated as equivalence.

\subsection*{Controls}
On a single earlier split of each pooled cohort, the cancer-type control was a logistic model whose only input was cancer type. The zero-pattern control used a binary indicator of whether each selected feature was zero or missing. Cohort identifiability was estimated by predicting cancer type from molecular values or from the zero pattern. Label-permutation nulls used 500 permutations per endpoint, either across all patients or within each cancer. In the cancer-held-out diagnostic, a linear molecular model was trained on all but one cancer and tested on the omitted one, with features preselected without labels in the pooled cohort.

\subsection*{Paired within-cancer comparison}
For each model and endpoint, we averaged each patient's test probability across repeats separately for the cancer-specific and the pooled model, computed the ROC AUC of each within every cancer, and took the difference. Uncertainty came from 2,000 bootstrap resamples of patients, stratified by outcome within each cancer and shared across models so that differences between model families are paired. The family comparison is the mean change of foundation models (TabFM and the TabPFN versions evaluated at that endpoint) minus the mean change of tree ensembles (random forest, CatBoost, XGBoost and LightGBM), with a 95\% percentile interval.

\subsection*{Reproducibility of this report}
All figures and tables were generated from saved predictions and metrics; no model was refitted to produce them and no curve or interval was entered by hand. The scripts are \path{paper/analysis/reusability/compute_reusability_statistics.py} and \path{paper/analysis/reusability/build_reusability_figures.py} for the main figures, \path{paper/analysis/build_manuscript_assets.py} for tables and Extended Data, and \path{paper/analysis/within_cancer_pooling_contrast.py}. We consulted TRIPOD+AI for reporting and PROBAST+AI for risk of bias and applicability, but did not complete formal checklists.\citep{tripodai2024,probastai2024}
